\chapter{Introduzione}

Il presente documento descrive l'analisi e la specifica dei requisiti per il progetto \textit{BugBoard26}, 
una piattaforma software per la gestione collaborativa di \textit{issue} e \textit{bug} durante il ciclo di sviluppo di un prodotto software.  

\bigskip

L'obiettivo principale del sistema è fornire un ambiente:
\begin{itemize}
    \item \textbf{semplice}, per facilitare la segnalazione e la gestione dei problemi;
    \item \textbf{sicuro}, per garantire l'integrità e la riservatezza dei dati;
    \item \textbf{efficiente}, per supportare le attività di classificazione, assegnazione e monitoraggio delle problematiche software.
\end{itemize}

\bigskip

Il sistema integra funzionalità avanzate come:
\begin{itemize}
    \item filtri dinamici di ricerca,
    \item allegati di immagini,
    \item esportazione dei dati in diversi formati,
    \item generazione di report mensili sull'attività del team.
\end{itemize}

\bigskip

Il documento include:
\begin{itemize}
    \item la specifica dettagliata dei requisiti software;
    \item la modellazione dei casi d'uso secondo il formalismo di \textit{Alistair Cockburn} \cite{cockburn-use-cases};
    \item la descrizione dell'interazione tra utenti e sistema, supportata da diagrammi e mock-up.
\end{itemize}

\bigskip

Questo testo costituisce una guida completa per \textbf{sviluppatori}, \textbf{project manager} e \textbf{stakeholder}, 
fornendo una visione chiara e strutturata delle funzionalità e delle logiche alla base del progetto \textit{BugBoard26}.